
\appendix
\section{Artifact Appendix}

\subsection*{Abstract}
Our open source artifact is available on \href{https://github.com/microsoft/sarathi-serve}{GitHub}. This repository contains our implementation of \sysname as well as the harnesses and scripts for
running and plotting the experiments described in this paper.

This repository originally started as a fork of the vLLM project. \sysname is a  lightweight high-performance research prototype and doesn't have complete feature parity with open-source vLLM. We have only retained the most critical features and adopted the codebase for faster research iterations.

\subsection*{Scope}

This artifact allows the readers to validate the claims made
in the \sysname paper (the figures) and provides a means to replicate the experiments described. The artifact
can be used to set up the necessary environment, execute the
main results, and perform microbenchmarks, thus providing a comprehensive understanding of the key claims in \sysname. %

\subsection*{Contents}

The repository is structured as follows, the primary source code for the system is contained in directory \textit{/sarathi}. The implementations for custom CUDA kernels are within the \textit{/csrc} directory. All the scripts to reproduce the experiments are in \textit{/osdi-experiments} and finally, the trace files used for the experiments are stored in \textit{/data}.


\subsection*{Hosting}
You can obtain our artifacts from GitHub:
\href{https://github.com/microsoft/sarathi-serve}{GitHub}. The main branch of the Github repository is actively updated, but we will maintain clear and accessible instructions about our artifacts in an easily identifiable
README file. All the detailed instructions and README files to reproduce the experiments in the OSDI paper are available in the branch \textit{osdi-sarathi-serve}. 


\subsection*{Requirements}

\sysname has been tested with CUDA 12.1 on A100 and A40 GPUs. The specific GPU SKUs on which the experiments were performed and the parallelism strategies used are clearly explained in the README corresponding to the figures in the artifact, for ease of reproducibility.



